\chapter{Results}

\section{Descriptive Statistics}

\subsection{Retrieval Outcomes Across Conditions}

{\centering\noindent Insert Table \ref{tab:retrieval-metrics} about here\par}

{\centering\noindent Insert Figure \ref{fig:coverage} about here\par}

The two representations separate on anchor coverage@5 and overlap on MRR@5
(Table~\ref{tab:retrieval-metrics}). On coverage@5 the six text-only and the six
enriched conditions occupy disjoint ranges, .038 apart. On MRR@5, enriched
BM25 at 500 tokens falls below two text-only conditions, hybrid retrieval at
200 and at 500 tokens.

The source table marks these pooled values as not comparable across question
categories, whose achievable ceilings differ.

\subsection{Bivariate Associations}

Table~\ref{tab:descriptives} (not pre-specified) summarizes 480 observations,
one question under one condition, and carries no significance markers, for the
reasons given in its note. Indexing representation correlated .27 with
coverage@5 and .25 with MRR@5. The largest coefficients with either outcome belong
to no design factor: the number of gold anchors correlated $-$.41 with coverage@5 and
$-$.36 with MRR@5.

\section{Retrieval Quality}

\subsection{H1: Chunk Size}

H1: Chunk size affects retrieval quality, measured as anchor coverage@5 and
MRR@5.

{\centering\noindent Insert Table \ref{tab:stage1-tests} about here\par}

Chunk size was tested in two families, once within text-only and once within
metadata-enriched, on the three pairs per family that differ in size alone:
dense, BM25, and hybrid at \thesisChunkSizeSmall{} against
\thesisChunkSizeLarge{} tokens (Table~\ref{tab:stage1-tests}). H1 was not
supported. No such pair reached significance in either family, on either
metric, in any scope, pooled or per category: every testable comparison
returned a Holm-adjusted \emph{p} of exactly 1.000. Chunk size also carried the
smallest coefficients in Table~\ref{tab:descriptives}, $-$.02 with coverage@5
and .00 with MRR@5.

\subsection{H2: Retrieval Strategy}

H2: Retrieval strategy affects retrieval quality.

Strategy was tested in the same two families as chunk size, so H2 is not an
independent second test; six pairs per family differ in strategy alone, three
at each chunk size. H2 was not supported within text-only on either metric, nor
within metadata-enriched on coverage@5. It was supported within
metadata-enriched on MRR@5, where BM25 at \thesisChunkSizeLarge{} tokens lost to
both dense and hybrid retrieval at the same size (Table~\ref{tab:stage1-tests}).

The MRR@5 result is pooled and reproduces in no question category. Three
strategy-only comparisons, each dense against hybrid on coverage@5, had every
question tied and could not be tested.

\subsection{H3: Indexing Representation}

H3: Enriching the indexed string with document metadata improves retrieval
quality.

All six Family 1 pairs differ in representation alone, holding chunk size and
retrieval strategy fixed.

The H3 verdict differs by retrieval strategy (Table~\ref{tab:stage1-tests}).
For dense and hybrid, at both chunk sizes and on both metrics, it was supported, the enriched condition
winning in each pair. For BM25 at \thesisChunkSizeSmall{} tokens it was
supported on MRR@5 and not on coverage@5: enrichment surfaced no additional
anchors there but reordered those already retrieved. For BM25 at
\thesisChunkSizeLarge{} tokens H3 was not supported on either metric.

Indexing representation is the only manipulated variable in
Table~\ref{tab:descriptives} carrying a non-trivial coefficient. H3 was tested pooled and its verdict is the
pooled one; the effect is not evenly spread across question categories, which
is examined below.

\subsection{Post Hoc: Where the H3 Effect Sits}

The breakdown by category was not pre-specified as a test of H3, so it is
descriptive and no hypothesis is accepted or rejected on it. The three
categories hold
\thesisNumQuestionsFactual{} factual, \thesisNumQuestionsThematic{} thematic, and
\thesisNumQuestionsComparative{} comparative questions, the
\thesisNumQuestionsScored{} retrieval-scored items between them, each block
corrected within itself per metric.

Among factual questions the effect was present: enrichment was significant in
the same four Family~1 pairs as the pooled coverage result, dense and hybrid at
both chunk sizes, on both metrics, enrichment winning in each
(Table~\ref{tab:stage1-tests}). Factual coverage@5 ran from .071 to .286
without enrichment and from .429 to .786 with it, the largest movement in the
study.

Thematic questions showed no detectable movement, both arms near floor:
coverage@5 ran from .033 to .100 without enrichment and from .033 to .133
with it, and no thematic comparison approached significance.

For comparative questions the pattern pointed the right way without surviving
correction. Three raw \emph{p} values fell below .05, exact permutation values
of 1/32 and 1/64 as the appendix describes, and Holm correction at family size six,
on 11 questions, removed each.

Families~2 and 3 were null in every category block on both metrics, so the
pooled MRR@5 result of Family~3 reproduces in no category.

\section{Answer Behavior}

\subsection{The Case Partition}

Every generation fell into one of three cases, fixed by retrieval before any
answer was produced: Case A, unanswerable from the corpus; Case B, with no gold
anchor in the top \thesisTopK{}; and Case C, with at least one. Case B was the
majority, at 61.7\%.

\subsection{Outcomes by Case}

{\centering\noindent Insert Table \ref{tab:case-outcome} about here\par}

{\centering\noindent Insert Figure \ref{fig:case-outcome} about here\par}

Outcomes are read per case (Table~\ref{tab:case-outcome}), and nothing in this
subsection was pre-specified. No Case A record can be grounded and on target,
since unanswerable items carry no transcript identifiers. In Case C, pooled over all \thesisNumConditions{}
conditions, \thesisQeightCaseCAnsweredGroundedCorrect{} of \thesisQeightCaseC{}
records (\thesisQeightCaseCConversionPct{}) converted.

No difference in outcomes between conditions was detectable. Over the six
text-only conditions, pooled, 65 of 120 pairwise comparisons were testable and none
survived Holm correction; the smallest raw \emph{p} among them was .219, and
every adjusted value was 1.000. Across all three families, pooled and by
category, 415 of 1152 comparisons
were testable, the other 737 having no discordant pair, and the smallest raw
\emph{p} was .125, adjusted to .750; .125 is the exact test's floor at four
discordant pairs, as Methods describes.

\subsection{Conversion on Matched Items}

The matched-subset comparison, which was not pre-specified, asks whether the
pooled difference between the representations reflects a change in the
population compared rather than in behavior within it. It covers the items in
Case C under both representations: 38 matched items, 76 observations,
from 13 distinct questions. Conversion rose from 24 of the 38 items (63.2\%) under
text-only to 29 of 38 (76.3\%) under metadata-enriched,
two items converting under text-only alone and seven under metadata-enriched
alone. Neither the question-level test Methods specifies (\emph{p} = .500)
nor the item-level test (\emph{p} = .180) is significant, so no
statistically supported claim that enrichment improves Case C conversion on
comparable items follows, and the Methods caveat applies.

\section{Failure Modes}

The coverage split, the manual review's grouping of generator behaviors, and
Check 3 were not pre-specified, and none involves a test.

\subsection{Coverage Split}

The coverage split breaks Case C conversion down by the share of an item's gold
anchors that was retrieved (Table~\ref{tab:coverage-split}). Its figures have
no reference values to be validated against and are reported as computed and
unchecked.

{\centering\noindent Insert Table \ref{tab:coverage-split} about here\par}

Case C requires only one anchor, so a comparative item holding one of its two
is in Case C although what was retrieved cannot support the comparison.

At partial coverage, thematic observations mostly converted and comparative
observations mostly abstained.

\subsection{Manual Faithfulness Review}

The review read \thesisNumReviewItems{} generations of the reported arm, from 25
questions (Table~\ref{tab:review-categories}). A proportion over the packet is a
rate within that sample, except where a stratum was taken whole.

{\centering\noindent Insert Table \ref{tab:review-categories} about here\par}

All eight P2 items are comparative and in Case C. The one abstention read as
over-conservative, on a comparative item under enriched BM25 at 200 tokens, had
both sides in its context and is outside P2.

Every record labeled ungrounded was read, because its stratum was taken whole:
\thesisQeightAnsweredUngrounded{} of the arm's \thesisQeightNumRecords{}
records (1.7\%), a census rather than a sample. Four read as faithful, their
labels traced to I1 twice and to I2 and I3 once each. Four, all in P1, should
have abstained: each answered the single unanswerable item that drew answers,
once under each BM25 condition. One, under enriched hybrid retrieval at 200
tokens, was unfaithful. The reading thus confirms a failure in 5 of the arm's
\thesisQeightNumRecords{} records (0.9\%), all among those labeled ungrounded,
which is not the arm's rate of unfaithful answers, since P1 also reaches records
under other labels. In the
other direction, one reviewed record labeled grounded and on target, under
BM25 at 200 tokens, attributes one bank's
remarks to a company that appears in none of its retrieved chunks and nowhere
in the corpus.
I4, unlike the census, rests on a sample: 17 off-target records drawn from the
grid's \thesisQeightAnsweredGroundedOfftarget{}, of which 12 read as faithful, a
proportion that is not extrapolated to all
\thesisQeightAnsweredGroundedOfftarget{}.

Check 3, an automatic diagnostic, flags a record labeled grounded and on target
when a substantive figure in it is supported only by chunks from outside the
question's own transcripts. Every Check 3 figure is a lower bound, not a
measurement: its population is Case C records so labeled, which hold only 2 of
the 12 P1 items. One is \texttt{ref\_comp\_10} under BM25 at 500 tokens,
whose answer gives Huntington's first-quarter capital figures as Fifth Third's;
the outcome label read it as grounded and on target, every figure being present
in the retrieved text, and Check 3 flagged it. Its counts on the matched subset,
given in the note to Table~\ref{tab:review-categories}, are evidence in neither
direction at that size, and untested.

\section{Robustness Check}

The robustness check, which was not pre-specified, asks whether the Stage 2
conclusions survive a generator that shares almost nothing with the reported
one, and covers one open-weights model family and no other.
\thesisFlashRobustnessModel{} \parencite{qwen3.8-flash-next} produced
\thesisFlashRobustnessNumRecords{} records over the same \thesisNumConditions{}
conditions, at temperature 0, on the same seed and the same prompt hash. The
two generators differ in architecture, parameter count, attention design, an
n-gram embedding table, and KV cache precision at once, the last forced by the
runtime, and the reported arm was not re-run to match; no result here is
attributed to any one difference. Stage 1 was not re-run, so every retrieval
result above serves this arm unchanged.

{\centering\noindent Insert Table \ref{tab:robustness} about here\par}

Outcome labels largely agreed between the reported arm and this one
(Table~\ref{tab:robustness}). The five unanswerable disagreements all come from
\texttt{ref\_unans\_01}, the
single unanswerable item that drew answers from the reported arm, across five
conditions.

The Case C and coverage-split results replicate on this maximally different
generator, on denominators Stage 1 fixes. On the matched subset, at question
level, the conservative test, three of the 13 questions were discordant,
\emph{p} = .250. That value is the floor at three discordant questions: all
three ran the same way and the test still could not reach significance.

No manual review was carried out on this arm, and neither the generator
behaviors nor the attribution check was computed for it. Both are readings of
the reported arm's answer text and do not transfer to answers another model
wrote, so the absence is deliberate rather than a gap.
